\begin{casebox}
\textbf{知识点小案例：future 等子任务会饿死同一个线程池。}
特例是 4 个 worker 各自执行一个父任务，父任务提交一个子任务后同步 \code{future.get()}。此时所有 worker 都在等，ready queue 里有 4 个子任务却没人执行。这不是锁顺序死锁，而是运行时调度协议错误。由此推出规律：任务运行时不能让 worker 在持有执行资源时盲等同池子任务；应使用 continuation、helping 或结构化 join。工程上要输出 task state、ready count、worker state 和 wait reason，否则“线程池卡住”无法区分死锁、饥饿和背压。
\end{casebox}
